% =============================================================================
% Project : THE_PHILOSOPHY_OF_VIABILITY
% Book    : Book I — Ontology of Viability
% Book ID : VS.B01
%
% File    : 2_MAINMATTER/PART_I_IDENTITY_AND_MANIFESTATION/CHAPTER_02_INVARIANCE/sections/section_03_defining_invariants.tex
%
% Section / Chapter:
%   Section 03 — Defining Invariants
%   Chapter 02 — Invariance
%
% Purpose :
%   To establish the analytical requirements for formulating and
%   evaluating claims that a determination or relation is preserved
%   under a specified transformation, while distinguishing the
%   formal specification of invariance from its ontological grounds.
%
% Conceptual Policy:
%   1. An invariance claim must identify what is claimed to be
%      preserved, under which transformation, within what domain,
%      and according to which criterion.
%   2. The subject of an invariance claim must be distinguished
%      from the existent to which it belongs.
%   3. The transformation and the domain of its applicability
%      must be specified without presupposing universal validity.
%   4. The respect of preservation must be distinguished from
%      the subject being examined.
%   5. Exact preservation, approximate correspondence, and
%      resemblance must not be conflated.
%   6. The admissibility of a transformation is a constraint
%      on the scope of a claim, not an independent ontological
%      component of invariance.
%   7. Formal invariance does not automatically establish
%      ontological invariance.
%   8. Analytical conditions specify and evaluate a claim;
%      they do not create the feature claimed to be preserved.
%   9. No invariant is presumed to hold across unrestricted
%      transformations or domains.
%
% =============================================================================

\section{Defining Invariants}
\label{sec:defining-invariants}

\subsection{The Problem of Defining an Invariant}

The assertion that something remains unchanged is not,
by itself, an adequate definition of invariance.

Such an assertion remains indeterminate unless it specifies
what is claimed to remain preserved, under which
transformation, within what domain, and according to
which criterion.

A determination may remain preserved under one
transformation while undergoing alteration under
another. A relation may remain unchanged within
a restricted domain but fail to remain so beyond
that domain. A configuration may preserve one
organizational feature while its other determinations
or relations undergo modification.

Consequently, invariance cannot be treated as an
unrestricted property of an existent, nor can it
be established through the undifferentiated description
of two configurations as unchanged.

An invariance claim is a claim of preservation under
specified conditions. Its meaning and validity depend
upon the identification of the feature concerned,
the transformation considered, the domain within
which the claim applies, and the criterion by which
preservation is assessed.

The task of defining invariants is therefore not
to introduce a universal mechanism of preservation.
It is to establish the analytical requirements
through which particular claims of invariance
can be formulated, distinguished, and evaluated.

These requirements must be sufficiently precise
to prevent the concept from being extended beyond
what a particular claim actually establishes.

\subsection{The Subject of an Invariance Claim}

Every claim of invariance identifies something
whose preservation is being examined. This is
the subject of the invariance claim.

The subject may be a determination, relation,
structural feature, quantitative magnitude,
qualitative feature, or another identifiable
aspect associated with an existent.

It must be distinguished from the existent as
a whole.

An existent may undergo substantial transformation
while a particular determination or relation
remains preserved. Conversely, the preservation
of one feature does not establish that the
existent as a whole remains unchanged.

The subject of an invariance claim therefore
identifies the particular feature under examination
without implying that this feature exhausts
the constitution of the existent or determines
its identity.

The distinction is also important when the
subject is relational. A claim concerning the
preservation of a relation must identify the
relation under consideration and the relevant
terms or relata involved in it. The preservation
of a relational form does not necessarily imply
the preservation of every determination of
the relata, nor does it establish that the
relata themselves remain numerically identical.

The subject must be sufficiently determinate
for the preservation claim to be evaluated.
An unspecified assertion that the existent,
its structure, or its organization remains
unchanged leaves open which feature is actually
being claimed as invariant.

Thus, identifying the subject establishes
what the claim concerns; it does not yet
establish that preservation occurs.

\subsection{Transformation as the Relational Context of Invariance}

Invariance is intelligible only in relation
to a transformation.

A transformation provides the context within
which the preservation or alteration of a
specified determination or relation can be
examined. Without a transformation, there
is no specified relation of preservation
to assess.

The transformation must be characterized
sufficiently to distinguish the relevant
initial and resulting configurations and
to identify the kind of change under
consideration.

A transformation may involve changes in
determinations, relations, composition,
structure, or other aspects of an existent.
It may also be restricted to a particular
class of changes.

Not every difference between two configurations
establishes that one has been transformed
into the other. The relation between the
configurations must itself be specified
where the claim depends upon a transformation
connecting them.

Likewise, not every transformation preserves
a feature. Preservation is a possible result
to be established, not an outcome that follows
from the mere specification of a transformation.

The same feature may remain preserved under
one transformation and undergo alteration
under another. Accordingly, an invariant
cannot be identified without reference to
the transformation relative to which its
preservation is claimed.

Transformation is therefore not an additional
property of the invariant. It specifies
the relation under which a preservation
claim is made.

\subsection{The Domain of Applicability}

A transformation is not necessarily applicable
to every possible configuration or condition.
The domain of an invariance claim specifies
the range within which the relevant
transformation and preservation criterion
are being considered.

A domain may be restricted by the configurations
included, the conditions under which the
transformation is defined, or the range
of cases to which the claim is intended
to apply.

For example, a structural relation may
remain preserved across a specified class
of transformations while undergoing
alteration under transformations outside
that class.

The preservation established within the
restricted domain does not automatically
extend beyond it.

The domain must therefore be distinguished
from the transformation itself. The
transformation specifies the relevant
change; the domain specifies the range
within which that transformation and
the corresponding preservation claim
are considered applicable.

A domain is an analytical delimitation
of a claim. Its specification does not,
by itself, establish an independent
ontological realm or introduce a new
kind of existent.

Nor does a restricted domain imply that
the claim is arbitrary. A domain may be
grounded in the actual conditions under
which a transformation occurs or in the
specified range of configurations to
which the claim applies.

The validity of the claim nevertheless
remains bounded by that domain unless
further grounds justify its extension.

\subsection{The Respect of Preservation}

The respect of preservation specifies
the regard in which a determination
or relation is claimed to remain unchanged.

It must be distinguished from the subject
of the claim.

The subject identifies what is being
examined. The respect specifies the
aspect under which its preservation
is assessed.

A configuration, for example, may retain
a particular structural arrangement while
its quantitative determinations change.
A relation may preserve its formal
organization while the properties of
its relata undergo alteration.

In such cases, the claim of preservation
is meaningful only if the relevant
respect is identified.

The respect must also be distinguished
from the criterion of preservation.
Identifying a respect establishes the
aspect in which sameness is being
examined; the criterion specifies
what counts as preservation in that
respect.

Without this distinction, a claim may
move ambiguously between different
features or different meanings of
sameness.

The preservation of a quantitative
determination, for example, cannot
be substituted without argument for
the preservation of a structural
relation. Likewise, the preservation
of a relational arrangement does
not establish the preservation of
the material constitution of the
terms participating in it.

An invariance claim must therefore
remain limited to the respect actually
specified and established.

\subsection{The Criterion of Preservation}

Identifying a subject, transformation,
domain, and respect does not yet establish
that preservation has occurred.

A criterion of preservation is required
to determine what must remain the same
for the specified subject to qualify
as invariant under the transformation
considered.

The criterion must correspond to the
nature of the subject and the respect
under examination.

Where the subject is a quantitative
determination, the criterion may require
the preservation of a specified value.
Where the subject is a structural relation,
the criterion may concern the maintenance
of a specified relational organization.
Where the subject is a configuration,
the criterion must identify which features
of that configuration are required to
remain preserved.

These examples do not establish a universal
criterion. The appropriate criterion
depends upon the subject and the claim
being formulated.

Several distinctions are necessary.

First, resemblance is not sufficient
to establish invariance. Two features
may appear similar without satisfying
the criterion required for preservation.

Second, correspondence is not necessarily
preservation. Two configurations may
correspond according to a mapping or
comparison rule while differing in
the respect relevant to the claim.

Third, approximate agreement must be
distinguished from exact preservation.
Where a criterion permits a tolerance
or degree of variation, that allowance
must be explicitly specified. A claim
of approximate invariance must not
be presented as exact invariance.

Fourth, persistence must not be substituted
for invariance. The continued occurrence
of a feature does not, by itself, establish
that the feature remains unchanged in
the respect under examination.

The criterion of preservation therefore
provides the standard against which
the claim is evaluated. It does not
create the preservation it is intended
to assess.

\subsection{The Scope and Admissibility of Transformation}

An invariance claim must specify which
transformations are included within
its scope.

Not every conceivable transformation
is relevant to every claim. A transformation
must be admissible relative to the
subject, domain, and respect under
consideration.

Admissibility concerns whether a
transformation belongs to the specified
class of transformations against which
the preservation claim is being assessed.

For example, a feature may be preserved
under a defined class of structural
transformations. This does not establish
that it remains preserved under changes
that fall outside that class or alter
the conditions required for the feature
to be identified in the first place.

The admissibility of a transformation
must therefore be justified by the
scope of the claim rather than assumed
from the mere possibility of describing
a change between two configurations.

This requirement does not introduce
a separate ontological component of
invariance. It clarifies the limits
of the transformation class relevant
to the claim.

The scope of validity is determined
by the combined specification of the
subject, transformation, domain, respect,
and preservation criterion, together
with the admissibility conditions
governing the transformations included.

A claim established under one set of
conditions cannot be generalized
automatically to a broader set.

\subsection{Formal Invariance and Ontological Invariance}

A distinction must be maintained between
formal invariance and ontological invariance.

A formal system may define an invariant
through a rule, operation, equation,
mapping, or transformation that preserves
a specified mathematical or logical
feature.

Such a result establishes preservation
relative to the formal structure and
the conditions defined within that
system.

An ontological claim, however, concerns
whether the feature in question is
actually preserved in the relevant
determinations or relations of an
existent.

The validity of a formal result and
the ontological status of the feature
represented by that result are therefore
analytically distinguishable.

A formal representation may describe,
model, or express an aspect of an
ontological relation. It may also
provide a rigorous means of investigating
whether a proposed feature remains
preserved under specified conditions.

Nevertheless, the formal preservation
of a mathematical or logical expression
does not automatically establish that
the corresponding feature is an
ontological invariant of an existent.

An ontological application requires
an account of the relation between
the formal structure and the existent
or feature to which the formal result
is being applied.

This requirement does not imply that
formal methods are ontologically
irrelevant. It establishes that their
results must not be extended beyond
their justified scope.

Formal invariance may contribute to
the investigation of ontological
invariance, but the transition from
a formal result to an ontological
claim requires an additional argument.

\subsection{The Analytical Conditions of an Invariance Claim}

The preceding analysis distinguishes
the principal elements required to
formulate and evaluate an invariance
claim.

These elements do not all have the
same logical function.

The subject, transformation, domain,
respect, and criterion specify the
content of the preservation claim.
The admissibility of transformations
and the scope of validity delimit
the conditions under which the claim
is applicable. Ontological grounding
concerns whether the result established
within an analytical or formal
framework warrants an ontological
interpretation.

They must therefore not be treated
as independent ontological components
or as successive stages through which
an invariant necessarily comes into
existence.

Their relations are summarized in
the following table.

% =============================================================================
% Project : THE_PHILOSOPHY_OF_VIABILITY
% Book    : Book I — Ontology of Viability
% Book ID : VS.B01
%
% File    : 4_ASSETS/tables/VS.B01.P01.C02.S03.SS10.P001.TL01.conditions_of_an_invariance_claim.tex
%
% Section / Chapter / Asset:
%   Chapter 02 — Invariance
%   Section 03 — Defining Invariants
%   Table TL01 — Conditions of an Invariance Claim
%
% Asset ID:
%   VS.B01.P01.C02.S03.SS10.P001.TL01.conditions_of_an_invariance_claim
%
% Purpose :
%   To distinguish the principal analytical elements and validity
%   conditions required for formulating and evaluating an invariance
%   claim.
%
% Conceptual Policy:
%   1. The table presents analytical requirements, not ontological
%      components.
%   2. Specification, scope, and ontological grounding have distinct
%      logical functions.
%   3. The table does not presuppose that an invariant exists.
%   4. The requirements are not presented as a universal formal
%      algorithm.
%
% =============================================================================

\begin{table}[htbp]
    \centering

    \caption{
        Analytical Conditions for Formulating and Evaluating
        an Invariance Claim
    }
    \label{tab:conditions-of-an-invariance-claim}

    \vspace{0.5em}

    \small

    \renewcommand{\arraystretch}{1.45}

    \begin{tabular}{
        |p{0.23\textwidth}
        |p{0.65\textwidth}|
    }
        \hline

        \textbf{Analytical Element}
        &
        \textbf{Question and Function}
        \\ \hline

        \textbf{Subject}
        &
        What determination, relation, structure, or other
        identifiable feature is claimed to be preserved?
        \\ \hline

        \textbf{Transformation}
        &
        Under which specified change or class of changes
        is preservation being examined?
        \\ \hline

        \textbf{Domain}
        &
        Within which range of configurations, conditions,
        or cases does the transformation and the claim apply?
        \\ \hline

        \textbf{Respect}
        &
        In what particular regard is the subject claimed
        to remain preserved?
        \\ \hline

        \textbf{Preservation Criterion}
        &
        What standard determines whether the specified
        subject remains unchanged in the relevant respect?
        \\ \hline

        \textbf{Transformation Admissibility}
        &
        Which transformations legitimately belong to the
        specified class against which preservation is assessed?
        \\ \hline

        \textbf{Scope of Validity}
        &
        What are the limits within which the preservation
        claim is justified, and beyond which it cannot
        be generalized without further argument?
        \\ \hline

        \textbf{Ontological Grounding}
        &
        What justifies interpreting the established
        preservation as a feature of the relevant existent
        or relation rather than merely as a formal result?
        \\ \hline

    \end{tabular}

    \normalsize

    \vspace{0.4em}

    \begin{minipage}{0.92\textwidth}
        \footnotesize
        \textit{Note.}
        The elements listed above have different analytical
        functions. They specify the content, applicability,
        limits, and ontological interpretation of an
        invariance claim. They are not independent
        ontological components and do not imply that
        an invariant necessarily exists.
    \end{minipage}

\end{table}

% =============================================================================
% END OF TABLE TL01
% =============================================================================

The table is an analytical guide
to the formulation of invariance
claims. It does not constitute an
exhaustive taxonomy of all possible
forms of invariance, nor does it
presuppose that every claim must
be expressed through an identical
formal procedure.

The requirements must be interpreted
in relation to the particular subject
and transformation under consideration.

A claim that identifies a subject
but leaves its transformation
unspecified remains incomplete.
A claim that specifies a transformation
but leaves its domain or respect
indeterminate does not establish
the full scope of preservation.
A claim that identifies a formal
invariant without justifying its
ontological application does not,
by that fact alone, establish an
ontological invariant.

The framework clarifies the conditions
of a meaningful claim without
converting those conditions into
independent features of the existent.

\subsection{The Conditional Character of Invariance}

The conditions established above
demonstrate that an invariance claim
is necessarily conditional.

This conditional character does not
mean that invariance is subjective
or that preservation depends upon
an arbitrary act of description.

It means that a claim of preservation
has determinate meaning only in
relation to the feature concerned,
the transformation considered,
the domain of applicability, and
the criterion by which preservation
is assessed.

The analytical specification of
these conditions does not create
the invariant. It establishes what
is being claimed and the circumstances
under which the claim may be justified.

The distinction between analytical
specification and ontological grounding
must therefore be maintained.

Analytical specification concerns
the formulation and evaluation of
the claim. Ontological grounding
concerns whether the feature identified
by that claim is genuinely preserved
in the relevant existent or relation.

Neither should be substituted for
the other.

An invariance claim may be precisely
formulated and nevertheless fail
because its proposed preservation
does not obtain. Conversely, a
feature may be preserved even when
an adequate account of that preservation
has not yet been formulated.

The distinction prevents the conditions
of description from being mistaken
for the conditions of existence.

\subsection{Synthesis}

An invariant cannot be adequately
defined as something that simply
remains unchanged.

A meaningful invariance claim requires
the identification of what is claimed
to be preserved, the transformation
under consideration, the domain
within which the claim applies,
the respect in which preservation
is assessed, and the criterion
by which preservation is established.

The admissibility of the transformation
and the scope of validity delimit
the applicability of the claim.
The distinction between formal
and ontological invariance determines
whether a formally established
preservation warrants an ontological
interpretation.

These conditions clarify the meaning
and limits of invariance without
presupposing that an invariant
necessarily exists.

They also prevent resemblance,
approximate correspondence, persistence,
continuity, and formal preservation
from being treated as interchangeable
with ontological invariance.

The framework developed here establishes
how claims of invariance are to be
specified and evaluated. It does
not establish that every existent
possesses an invariant, that every
transformation preserves something,
or that preservation is unlimited
across domains and respects.

The next section examines the relation
between invariance and transformation
more closely, particularly the possible
coexistence of preservation and
alteration within a single transformation.

% =============================================================================
% END OF SECTION 03 OF CHAPTER 02
% =============================================================================